\scanpage{173}
\mspage{47}

\uline{Dém} Soit $(\deleted{\unsure{X}}Z_\alpha)_{\alpha\in A}$ une famille de sous-objets
de $X$ tels que $\deleted{\unsure{X}}Z_\alpha\in\mathrm{Filt}^{\unsure{\pi|\varpi}}$, et $\mathrm{card}\,A\leq\pi$.
Soit $Z=\coprod_{\alpha\in A}\underset{\unsure{,,}}{\deleted{\unsure{X}}Z_\alpha}$, on a $Z\in\mathrm{Filt}^{\pi}$ par d).
Le plus petit sous-objet de $X$ majorant les
$Z_\alpha$ n'est autre que $\mathrm{Im}(Z\to X)$, qui
est dans $\mathrm{Filt}^{\pi}$ par $\mathrm d_{\unsure{2|1}}$. Prouvons que
$\underset{\alpha\in\mathrm{Ss}_{\pi}(X)}{\underrightarrow{\mathrm{lim}}}\,\deleted{\unsure{X}}Z_{\alpha}\overset{\sim}{\longrightarrow}X$ est iso. \deleted{On sait qu'} \inserted{Quand on sait que c'est un mono}
\unsure{--|a} $X=\underset{I}{\underrightarrow{\mathrm{lim}}}\,X_i$, où $(X_i)_{i\in\unsure{J|I}}$ est
un système inductif filtrant, $I$ un
ordonné grand devant $\pi$, \inserted{les $X_i$ dans $\mathrm{Filt}^{\pi}$.} Soit, pour tout
$i\in I$, $X'_i$ l'image de $X_i$ dans $X$.
\inserted{On a donc $X'_i\in\mathrm{Filt}^{\pi}$ par $\mathrm d_2$).}
Utilisant le fait que les $X_i\to X'_i$ sont
épi, on voit de suite que $X=\underset{I}{\underrightarrow{\mathrm{lim}}}\,X'_i$.
Considérons l'application
\diagram{\begin{tikzcd}[ampersand replacement=\&,row sep=0.2em,column sep=2em]
 I \arrow[r] \& \mathrm{Ss}_{\pi}(X) \\
 i \arrow[r,mapsto] \& X'_i
\end{tikzcd}}
C'est une application croissante, \inserted{soit $J$ son image,
partie filtrante grande devant $\pi$ de $\mathrm{Ss}_{\pi}(X)$} \deleted{pour que l'} telle
\deleted{est cofinale. Comme $I,\unsure{I'|J'}$ sont des ens. ordonnés
filtrants, \unsure{tout revient} cela signifie que pour
tout $Z_\alpha$ dans $\mathrm{Ss}_{\pi}(X)$, existe un $i$ tel que
$Z_\alpha\subset X'_i$.} que $\underset{\alpha\in J}{\underrightarrow{\mathrm{lim}}}\,Z_\alpha\overset{\sim}{\longrightarrow}X$.

% archive sheet number: 172
